\documentclass[12pt,a4paper]{article}

\input{preamble}

\course{Algebra I}
\sheetnumber{3}
\lecturer{Lecturer}
\duedate{15 October 2026}

\showdraftnotestrue

\begin{document}

\makesheetheader

\begin{exercise}
    \label{exercise:basic}
    Let \(V\) be a finite-dimensional vector space.
    \begin{enumerate}
        \item Define a canonical map \(V \to V^{\vee\vee}\).
            \label{item:dual-dual-map}
        \item Prove that this map is an isomorphism.
            \label{item:dual-dual-isomorphism}
    \end{enumerate}
\end{exercise}

\showsolutionstrue
\begin{solution}
    For \(v \in V\), define
    \[
        \Phi(v) \colon V^\vee \to k,
        \qquad
        \Phi(v)(\lambda)=\lambda(v).
    \]
    This gives a linear map \(\Phi \colon V \to V^{\vee\vee}\), which is the map
    asked for in \Cref{item:dual-dual-map}. After choosing a basis of \(V\), the
    map \(\Phi\) is identified with the usual identification of a finite-dimensional
    vector space with its double dual. Hence \(\Phi\) is an isomorphism, proving
    \Cref{item:dual-dual-isomorphism}.
\end{solution}
\showsolutionsfalse

\draftnote{Hide the previous solution in the future.}

\begin{exercise}
    \label{exercise:exact-sequence}
    Let
    \[
        0 \to A^{\vee} \to B^{\vee} \to C^{\vee} \to 0
    \]
    be a short exact sequence of finite-dimensional vector spaces. Use
    \Cref{item:dual-dual-isomorphism} of \Cref{exercise:basic} to describe the
    dual exact sequence.
\end{exercise}

\begin{solution}
    Dualizing gives an exact sequence
    \[
        0 \to C^{\vee\vee} \to B^{\vee\vee} \to A^{\vee\vee} \to 0.
    \]
    The finite-dimensionality assumption allows us to identify each vector space
    with its double dual, using \Cref{item:dual-dual-isomorphism} of
    \Cref{exercise:basic}.
    So we obtain a short exact sequence
    \[
        0 \to C \to B \to A \to 0.
    \]
\end{solution}

\end{document}
